\documentclass{article}
\usepackage{amsmath}
\usepackage{graphicx}
\usepackage{amssymb}

\title{Lecture 15: Plane wave propagation}
\author{David Al-Attar, Michaelmas term 2024}
\date{}

\begin{document}

\maketitle

\section*{Outline and motivation}

In this lecture, we step back from realistic applications that require numerical calculations to consider the propagation of plane waves with within an elastic whole space with constant material parameters. We show that under reasonable restrictions on the elastic tensor, up to three different types of plane wave can propagate in each direction. The existence of multiple wave types is a fundamental property of elasticity, and accounts for much of the complexity of the observed seismic wavefield. In the case of isotropic materials a simple closed-form solution of the problem can be found, and in this manner we arrive at P- and S-waves. We then apply perturbation theory to understand how these solutions are modified by the slight anisotropy present in some regions of the Earth. To conclude we discuss briefly the interaction of plane waves with a material interface, showing that reflections, transmissions and conversions can occur.

\section*{Derivation of the Christoffel equation}

In a previous lecture we obtained the linearised equations of motion
\begin{equation}
\rho\frac{\partial^{2}u_{i}}{\partial t^{2}}-\frac{\partial}{\partial x_{j}}(A_{ijkl}\frac{\partial u_{k}}{\partial x_{l}})=0,
\end{equation}
where $\rho$ is the referential density, and $A_{ijkl}$ the elastic tensor. We recall that the elastic tensor possesses the symmetry $A_{ijkl}=A_{klij}$, and will start to see the significance of this property as we proceed today. In the first part of this lecture we will assume that the reference body is equal to $\mathbb{R}^{3}$, and that the material parameters are independent of position. We say that a body is homogeneous if its material parameters are independent of position, and otherwise it is said to be heterogeneous. Note that though the terms "homogeneous" and "isotropic" are sometimes confused but they mean completely different things. The first refers to the invariance of properties with spatial position, while the latter is an invariance of the properties with respect to orientation at a fixed point. A homogeneous medium could be anisotropic, while a heterogeneous one can be isotropic.

We begin by looking for plane wave solutions of eq.(1) taking the form
\begin{equation}
u_{i}(x,t)=f(t-p_{j}x_{j})a_{i}
\end{equation}
where p is the real-valued slowness vector, a is the polarisation vector, and f is an arbitrary function. The fact that the shape of the plane wave does not change during propagation shows that, if such waves exist, they are non-dispersive. The above displacement vector is seen to be constant on the wavefronts
\begin{equation}
t-p_{j}x_{j}=\text{constant},
\end{equation}
and these move through space in the direction of p at the phase speed
\begin{equation}
c=\frac{1}{||p||}
\end{equation}
Here we see that the slowness vector is used to specify both the speed and direction of propagation of the plane wave, while the polarisation vector gives the orientation of the associated particle motion. To see if such plane wave solutions are possible, we simply substitute the displacement vector in eq. (2) into the equations of motion. First, we find readily that
\begin{equation}
\frac{\partial^{2}u_{i}}{\partial t^{2}}=f''(t-p_{j}x_{j})a_{i},
\end{equation}
where $f''$ denotes the second derivative of f with respect to its argument. To calculate the remaining terms in the equations of motion, we first note that
\begin{equation}
\frac{\partial}{\partial x_{j}}(A_{ijkl}\frac{\partial u_{k}}{\partial x_{l}})=A_{ijkl}\frac{\partial^{2}u_{k}}{\partial x_{j}\partial x_{l}}
\end{equation}
because the elastic tensor is constant. From the definition of u above, we can use the chain rule to show that
\begin{equation}
\frac{\partial u_{k}}{\partial x_{l}}=-f'(t-p_{m}x_{m})a_{k}p_{l},
\end{equation}
and hence that
\begin{equation}
\frac{\partial^{2}u_{k}}{\partial x_{j}\partial x_{l}}=f''(t-p_{m}x_{m})a_{k}p_{j}p_{l}.
\end{equation}
Combining these calculations, we arrive at
\begin{equation}
\frac{\partial}{\partial x_{j}}(A_{ijkl}\frac{\partial u_{k}}{\partial x_{l}})=f''(t-p_{m}x_{m})A_{ijkl}p_{j}p_{l}a_{k}
\end{equation}
Cancelling the common factor of $f''$, it follows that the equations of motion will be satisfied if and only if
\begin{equation}
(\rho\delta_{ik}-A_{ijkl}p_{j}p_{l})a_{k}=0.
\end{equation}
To simplify this result, it will be useful to define the Christoffel matrix $\Gamma(p)$ for a given slowness vector p to have components
\begin{equation}
\Gamma_{ik}(p)=\frac{1}{\rho}A_{ijkl}p_{j}p_{l}
\end{equation}
Using this definition, then eq. (10) can be written using matrix-vector notation as
\begin{equation}
[\Gamma(p)-I]a=0
\end{equation}
which is known as the Christoffel equation. This equation acts as a restriction on the possible values of the slowness and polarisation vectors for which plane elastic wave propagation is possible. As a first observation, we note that this is a homogeneous system of linear equations for a. It always has the trivial solution $a=0$, while non-trivial solutions will only exist if
\begin{equation}
\det[\Gamma(p)-I]=0
\end{equation}
This is an algebraic condition on the slowness vector alone, and the totality of its solutions defines what is known as the slowness surface.

\section*{Reduction to a symmetric eigenvalue problem}

To obtain solutions of the Christoffel equation it is useful to first note two properties of the Christoffel matrix. The first is that
\begin{equation}
\Gamma(\lambda p)=\lambda^{2}\Gamma(p),
\end{equation}
for any $\lambda$ and p. We know that the slowness vector sets both the direction and speed of the plane wave. It is, therefore, useful to write it in the form
\begin{equation}
p=\frac{1}{c}\hat{p}
\end{equation}
where c is the phase speed, and $\hat{p}$ a unit vector. Having done this, eq. (14) allows us to rewrite the Christoffel equation as
\begin{equation}
\Gamma(\hat{p})a=c^{2}a
\end{equation}
If we focus on plane waves in a fixed direction, then we have arrived at an eigenvalue problem with $c^{2}$ the eigenvalue, and a the eigenvector. Next, for any p the Christoffel matrix is symmetric. Indeed, by definition we have
\begin{equation}
\Gamma_{ki}(p)=\frac{1}{\rho}A_{kjil}p_{j}p_{l}=\frac{1}{\rho}A_{ilkj}p_{j}p_{l}=\frac{1}{\rho}A_{ijkl}p_{j}p_{l}=\Gamma_{ik}(p).
\end{equation}
where in the first equality we have used the hyperelastic symmetry, while the second has simply relabelled the dummy indices. This result is very useful, because we can now see that eq.(16) is an eigenvalue problem for a symmetric matrix, and hence of a form that admits simple solutions.

\subsection*{Qualitative solution of the eigenvalue problem}

With these facts at hand we can describe in qualitative terms the plane wave solutions possible in a fixed propagation direction. Eq.(16) is a symmetric eigenvalue problem in $\mathbb{R}^{3}$, and so has three eigenvalues $c_{1}^{2}$, $c_{2}^{2}$, $c_{3}^{2}$, with corresponding eigenvectors $a_{1}$, $a_{2}$, $a_{3}$. Assuming that the squared phase speeds are positive, this gives six possible phase speeds $\pm c_1, \pm c_{2}, \pm c_{3}$, with the pairs corresponding to waves propagating either forward and backwards along $\hat{p}$. Moreover, eigenvectors belonging to distinct eigenvalues are necessarily orthogonal. What does this mean physically? In a given propagation direction, $\hat{p}$, there are up to three different plane wave solutions possible (or six if we include those propagating backwards). The phase speed of each wave is set by the corresponding eigenvalue of the Christoffel matrix, and its polarisation vector is determined by the associated eigenvector. Finally, the polarisation vectors for plane waves with different phase speeds are necessarily orthogonal.

Note that the Christoffel matrix might have degenerate eigenvalues, meaning that an eigenvalue has more than one linearly independent eigenvector. If this occurs, then these eigenvectors need not be orthogonal, but we can always pick an orthogonal basis that spans the eigenspace. We will see below that such degeneracies occur in the case of isotropic materials, and hence this situation cannot be discounted.

The above discussion has implicitly assumed that the eigenvalues of the Christoffel matrix are all non-negative, and hence lead to well-defined phase speeds. This requirement places a physical constraint on the components of the elastic tensor. We will not look into this issue in any detail, but note that it is met in all real materials. Moreover, in a solid the three phase speeds are always positive, while in a fluid there is one positive phase speed and two that vanish.

\begin{figure}[h]
\centering
\includegraphics[width=0.8\textwidth]{placeholder_image.png}
\caption{Plane sections through the slowness surface associated with two randomly selected anisotropic elastic materials. The three sheets of the slowness surface are shown in different colours.}
\label{fig:slowness}
\end{figure}

\subsection*{Determination of the slowness surface}

Having now determined the form of solutions of the Christoffel equation in a fixed propagation direction all we need do is vary $\hat{p}$ over the unit sphere to obtain a complete description of the possible plane waves in the body. For each direction we get (allowing for possible degeneracies) three possible phase speeds, and hence three values of p. Let us write $c_{1}(\hat{p})$, $c_{2}(\hat{p})$, $c_{3}(\hat{p})$ for the three non-negative phase speeds along $\hat{p}$ and order them such that
\begin{equation}
c_{1}(\hat{p})\le c_{2}(\hat{p})\le c_{3}(\hat{p}),
\end{equation}
where we have to use $\le$ to allow for possible degeneracies. The corresponding values of the slowness vector are
\begin{equation}
p=\frac{1}{c_{1}(\hat{p})}\hat{p}, \quad p=\frac{1}{c_{2}(\hat{p})}\hat{p}, \quad p=\frac{1}{c_{3}(\hat{p})}\hat{p}.
\end{equation}
As $\hat{p}$ varies over the unit sphere, these slowness vectors build up the slowness surface. This surface is obviously closed, and is symmetric under reflection about the origin. If there are no degeneracies and all phase speeds are positive, then it comprises three distinct sheets. If degeneracies occur along certain propagation directions, then two or more sheets will touch. We note also that the defining relation in eq.(13) can now be written
\begin{equation}
[||p||^{2}c_{1}(\hat{p})^{2}-1][||p||^{2}c_{2}(\hat{p})^{2}-1][||p||^{2}c_{3}(\hat{p})^{2}-1]=0,
\end{equation}
with each factor defining a sheet of the slowness surface.

For a general elastic tensor the determination of the slowness surface requires solution of the symmetric eigenvalue problem in each possible propagation direction. Such calculations are, of course, very easy to do numerically. In fig.1 we illustrate these properties of the slowness surface by plotting slices through a number of examples that have been calculated numerically. We note, in particular, that the innermost sheet of the slowness surface is always convex in these examples. This is, in fact, a general property, and a proof might be sought if someone is particularly interested\footnote{As a hint, note that eq. (13) is a sextic polynomial in p, and hence along a fixed line in $\mathbb{R}^{3}$ there can be no more than six roots.}.

\section*{The Christoffel equation in an isotropic material}

In the case of an isotropic material a simple closed-form solution of the Christoffel equation can be obtained, and in this manner we arrive at P- and S-waves that you have likely seen previously. We recall from the last lecture that in an isotropic material the elastic tensor takes the form
\begin{equation}
A_{ijkl}=\lambda\delta_{ij}\delta_{kl}+\mu(\delta_{ik}\delta_{jl}+\delta_{il}\delta_{jk}),
\end{equation}
with $\lambda$ and $\mu$ the Lamé parameters. A simple calculation then yields
\begin{equation}
\Gamma_{ik}(\hat{p})=\frac{\lambda+\mu}{\rho}\hat{p}_{i}\hat{p}_{k}+\frac{\mu}{\rho}\delta_{ik}.
\end{equation}
At this stage we could write out $\Gamma_{ik}(\hat{p})$ explicitly as a matrix and calculate its eigenvalues and eigenvectors in the usual manner. There is, however, a quicker way to proceed which is worth learning. First we rewrite the expression as
\begin{equation}
\Gamma_{ik}(\hat{p})=\frac{\lambda+2\mu}{\rho}\hat{p}_{i}\hat{p}_{k}+\frac{\mu}{\rho}(\delta_{ik}-\hat{p}_{i}\hat{p}_{k}).
\end{equation}
We can now identify two terms with a geometric significance. $\hat{p}_{i}\hat{p}_{k}$ are the components of a matrix that projects a vector parallel to $\hat{p}$ while $\delta_{ik}-\hat{p}_{i}\hat{p}_{k}$ are those of one that projects orthogonally to $\hat{p}$. It follows that if we take the polarisation vector parallel to the propagation direction we obtain
\begin{equation}
\Gamma(\hat{p})\hat{a}=\frac{\lambda+2\mu}{\rho}\hat{a}.
\end{equation}
This shows that $a\propto \hat{p}$ is an eigenvector of the Christoffel matrix with phase speed
\begin{equation}
\alpha=\sqrt{\frac{\lambda+2\mu}{\rho}}.
\end{equation}
Next, if we let the polarisation vector be orthogonal to $\hat{p}$ we find
\begin{equation}
\Gamma(\hat{p})\hat{a}=\frac{\mu}{\rho}\hat{a}
\end{equation}
and hence have found two further linearly independent eigenvectors with a common eigenvalue whose associated phase speed is
\begin{equation}
\beta=\sqrt{\frac{\mu}{\rho}}.
\end{equation}
Physically, we have now shown that there are two types of plane waves possible in an isotropic whole space. The first are P-waves which have phase speed $\alpha$ independent of $\hat{p}$, and whose polarisation vector lies parallel to the propagation direction. The next are S-waves which have phase speed $\beta$ which is again independent of $\hat{p}$, and whose polarisation vectors must lie in the plane orthogonal to the propagation direction. It can be shown that in any real isotropic solid
\begin{equation}
\mu>0 \quad \kappa=\lambda+\frac{2}{3}\mu>0
\end{equation}
and hence we always have $\alpha>\beta>0$. The slowness surface, therefore, has two spherical sheets, with the inner one corresponding to P-waves, while the outer one for S-waves is two-fold degenerate. Note finally that in the case of an elastic fluid the above solution still applies, but as $\mu=0$, there are no S-waves that can propagate.

\section*{Perturbation theory for slightly anisotropic materials}

We noted in the last lecture that seismic waves generally see the Earth as being close to isotropic. It is, therefore, of interest to ask how the plane wave solutions just obtained in the isotropic case are modified by the presence of slight anisotropy. We recall that in considering propagation in direction $\hat{p}$ we have to solve the symmetric eigenvalue problem
\begin{equation}
\Gamma(\hat{p})a=c^{2}a
\end{equation}
Let us suppose that the elastic tensor takes the following form
\begin{equation}
A_{ijkl}=\lambda\delta_{ij}\delta_{kl}+\mu(\delta_{ik}\delta_{jl}+\delta_{il}\delta_{jk})+s~\delta A_{ijkl}
\end{equation}
Here s is a perturbation parameter, with $s=0$ corresponding to the isotropic case, and $\delta A_{ijkl}$ an anisotropic perturbation. The corresponding Christoffel matrix can then be written
\begin{equation}
\Gamma_{ik}(\hat{p})=\frac{\lambda+\mu}{\rho}\hat{p}_{i}\hat{p}_{k}+\frac{\mu}{\rho}\delta_{ik}+s~\delta\Gamma_{ik}(\hat{p}),
\end{equation}
where the zeroth-order part is that for an isotropic material given in eq.(22), while the first order term is given by
\begin{equation}
\delta\Gamma_{ik}(\hat{p})=\frac{1}{\rho}\delta A_{ijkl}\hat{p}_{j}\hat{p}_{l}.
\end{equation}
The eigenvalue problem to be solved can, therefore, be written
\begin{equation}
[\Gamma(\hat{p})+s~\delta\Gamma(\hat{p})]a=c^{2}a
\end{equation}
where $\Gamma(\hat{p})$ is here understood to be the isotropic part of the Christoffel matrix. While a numerical solution of the problem is always possible, significant insight can be gained through an application of first-order eigenvalue perturbation theory. Some of you will have seen these ideas before (for example, in quantum mechanics), but we will derive the necessary results from first-principles. In the isotropic case we know that the possible solutions correspond to P- and S-waves, with the latter being two-fold degenerate. For P-waves we can apply non-degenerate perturbation theory. This assumes that the phase speed takes the form
\begin{equation}
c=\alpha+s~\delta\alpha+\cdot\cdot\cdot,
\end{equation}
while the polarisation vector can be written
\begin{equation}
a=\hat{p}+s\delta a+\cdot\cdot\cdot
\end{equation}
Putting these expansions into the eigenvalue problem and isolating the first-order terms in s, we find that
\begin{equation}
\Gamma(\hat{p})\delta a+\delta\Gamma(\hat{p})\hat{p}=\alpha^{2}\delta a+2\alpha\delta\alpha\hat{p}.
\end{equation}
This equation involves both $\delta\alpha$ and $\delta a$, but we can eliminate the latter by taking the inner product with $\hat{p}$ and noting that
\begin{equation}
\hat{p}\cdot[\Gamma(\hat{p})\delta a]=[\Gamma(\hat{p})\hat{p}]\cdot\delta a=\alpha^{2}\hat{p}\cdot\delta a
\end{equation}
where we have used the symmetry of the Christoffel equation along with the zeroth-order eigenvalue problem. In this manner we arrive at the simple expression
\begin{equation}
\delta\alpha=\frac{1}{2\alpha}\hat{p}\cdot\delta\Gamma(\hat{p})\hat{p}
\end{equation}
for the first-order change to the phase speed. The first-order perturbation to the polarisation vector could then also be determined, but we will not work through the details. What is physically important is that in a slightly anisotropic material we have shown that there will be a plane wave that has phase speed close to that of a P-wave, and whose polarisation vector is almost aligned with the propagation direction. Note that, however, the phase speed will now generally depend on the propagation direction as determined by eq. (38). Such a plane wave is known as a quasi P-wave.

The S-wave case can be handled similarly, but is complicated by the unperturbed eigenspace being degenerate. Let $\hat{t}_{1}$ and $\hat{t}_{2}$ be two unit vectors orthogonal to the propagation direction. In the absence of anisotropy we know that the polarisation vector can be any vector orthogonal to $\hat{p}$. It therefore seems reasonable to look for solutions
\begin{align}
c&=\beta+s~\delta\beta+\cdot\cdot\cdot, \\
a&=a_{1}\hat{t}_{1}+a_{2}\hat{t}_{2}+s~\delta a+\cdot\cdot\cdot,
\end{align}
where $a_{1}$ and $a_{2}$ are zeroth-order terms to be determined. Putting these expansions into the eigenvalue problem, the first-order term is given by
\begin{equation}
\Gamma(\hat{p})\delta a+\delta\Gamma(\hat{p})(a_{1}\hat{t}_{1}+a_{2}\hat{t}_{2})=\beta^{2}\delta a+2\beta\delta\beta(a_{1}\hat{t}_{1}+a_{2}\hat{t}_{2}).
\end{equation}
In a similar manner we can eliminate $\delta a$ by taking the inner product with either $\hat{t}_{1}$ or $\hat{t}_{2}$, and doing this leads to the following pair of equations
\begin{align}
\hat{t}_{1}\cdot\delta\Gamma(\hat{p})(a_{1}\hat{t}_{1}+a_{2}\hat{t}_{2})&=2\beta\delta\beta a_{1}, \\
\hat{t}_{2}\cdot\delta\Gamma(\hat{p})(a_{1}\hat{t}_{1}+a_{2}\hat{t}_{2})&=2\beta\delta\beta a_{2}.
\end{align}
Written in matrix form, we have obtained
\begin{equation}
\begin{pmatrix} \hat{t}_{1}\cdot\delta\Gamma(\hat{p})\hat{t}_{1} & \hat{t}_{1}\cdot\delta\Gamma(\hat{p})\hat{t}_{2} \\ \hat{t}_{2}\cdot\delta\Gamma(\hat{p})\hat{t}_{1} & \hat{t}_{2}\cdot\delta\Gamma(\hat{p})\hat{t}_{2} \end{pmatrix} \begin{pmatrix} a_{1} \\ a_{2} \end{pmatrix} = 2\beta\delta\beta \begin{pmatrix} a_{1} \\ a_{2} \end{pmatrix}.
\end{equation}
This is a $2\times2$ symmetric eigenvalue problem whose solution determines both the zeroth-order polarisation vectors and the first-order changes in phase speed. It would not be difficult to solve in closed-form, but the qualitative behaviour is more important. In the unperturbed case S-waves are associated with a two-fold degenerate sheet of the slowness surface. For an anisotropic perturbation, the above eigenvalue problem will almost always have two distinct eigenvalues, and hence the S-wave degeneracy will be split. There will then be two distinct plane waves whose phase speeds depend on the propagation direction. They will, moreover, have orthogonal polarisation vectors lying almost orthogonal to the propagation direction; note the first-order term $\delta a$ will in general have a component along $\hat{p}$. These planes waves are known as quasi S-waves, and the process we have described is known as shear wave splitting. A similar phenomena occurs for waves in the Earth, and the observation of shear wave splitting provides an important constraint on anisotropy. The reason this is interesting is as follows. Most of the Earth appears to seismic waves to be isotropic. This is because anisotropic grains are generally oriented randomly, and hence the structure seen by seismic waves averages out to be isotropic. In some parts of the mantle, however, the seismic waves do see a slight anisotropy because there is a large-scale statistical alignment of grains due to convective flow. Thus, by making observations of shear wave splitting it has been possible to make inferences about the orientation of flow within parts of the mantle.

\section*{Reflection, transmission, and conversion at planar boundaries}

Suppose now that we are not dealing with a homogeneous whole space, but one comprised of two distinct homogeneous half spaces. Without loss of generality, we can put the interface at $x_{3}=0$. Within each half space we know that plane wave solutions of the equations of motion are possible. But to obtain a full solution of the problem we need to consider continuity conditions at the interface. Restricting attention to solids, these conditions are that the displacement $u_{i}$ and the traction $t_{i}=T_{ij}\hat{n}_{j}$ are continuous, which we write mathematically as
\begin{equation}
[u_{i}]_{-}^{+}=0, \quad [t_{i}]_{-}^{+}=0.
\end{equation}
To proceed it is useful to note that a general plane wave can be written using Fourier transforms as
\begin{equation}
u_{i}(x,t)=f(t-p_{j}x_{j})a_{i}=\frac{1}{2\pi}\int_{\mathbb{R}}\hat{f}(\omega)e^{i\omega(t-p_{j}x_{j})}d\omega~a_{i}
\end{equation}
and by linearity of the problem it is sufficient to consider harmonic plane waves
\begin{equation}
u_{i}(x,t)=e^{i\omega(t-p_{j}x_{j})}a_{i}.
\end{equation}
Within the lower half-plane let us consider such a plane wave incident on the boundary. The slowness, $p_{i}$, and polarisation vector, $a_{i}$, for this wave are fixed by the Christoffel equation within the lower medium, while to be upwards propagating we require $p_{3}>0$. To satisfy the equations of motion and continuity conditions we would expect some number of reflected and transmitted plane waves, and hence we assume that the displacement vector takes the form
\begin{equation}
u_{i}(x,t)=\begin{cases} e^{i\omega(t-p_{j}x_{j})}a_{i}+\sum_{m}e^{i\omega(t-p_{jm}^{-}x_{j})}a_{im}^{-} & x_{3}\le0 \\ \sum_{m}e^{i\omega(t-p_{jm}^{+}x_{j})}a_{im}^{+} & x_{3}\ge0 \end{cases}
\end{equation}
The first term in the expression for $x_{3}\le0$ denotes the incident wave whose slowness satisfies $p_{3}>0$. The next term then represents a sum over some number of reflected waves indexed by m, and with a minus subscript to emphasise that these terms are associated with the lower layer. Finally, the solution for $x_{3}\ge0$ is a similar sum over some number of transmitted phases. In order for the equations of motion to be satisfied, the slowness and polarisation vectors for each plane wave must solve the appropriate Christoffel equation.

\begin{figure}[h]
\centering
\includegraphics[width=\textwidth]{placeholder_image_2.png}
\caption{Graphical representations of Snell's law for the reflection and transmission of plane waves between two homogeneous elastic half-spaces. In each figure lower and upper plane sections of the slowness surface for the materials are shown. The slowness vector of the incident wave is plotted as a yellow arrow, and the continuation of the appropriate sheet of the slowness surface in the lower medium into the upper half of the plot has been included in yellow for convenience. Snell's law fixes the horizontal slowness of both the reflected and transmitted waves, and the intersections of the plotted vertical line with the two slowness surfaces gives the real slowness vectors of the reflected and transmitted waves. In the figure on the right, the incident horizontal slowness is such that evanescent reflected and transmitted waves must be included to satisfy the continuity of displacement and traction at the boundary.}
\label{fig:snell}
\end{figure}

(defined using the material parameter relevant to the half space considered). In particular, the slowness vectors must lie somewhere on the appropriate slowness surface. Continuity of displacement at $x_{3}=0$ implies
\begin{equation}
e^{i\omega(t-p_{1}x_{1}-p_{2}x_{2})}a_{i}+\sum_{m}e^{i\omega(t-p_{1m}^{-}x_{1}-p_{2m}^{-}x_{2})}a_{im}^{-}=\sum_{m}e^{i\omega(t-p_{1m}^{+}x_{1}-p_{2m}^{+}x_{2})}a_{im}^{+};
\end{equation}
while that for the traction is given by
\begin{multline}
-i\omega p_{l}A_{i3kl}^{-}e^{i\omega(t-p_{1}x_{1}-p_{2}x_{2})}a_{k}-\sum_{m}i\omega p_{lm}^{-}A_{i3kl}^{-}e^{i\omega(t-p_{1m}^{-}x_{1}-p_{2m}^{-}x_{2})}a_{km}^{-} \\
=-\sum_{m}i\omega p_{lm}^{+}A_{i3kl}^{+}e^{i\omega(t-p_{1m}^{+}x_{1}-p_{2m}^{+}x_{2})}a_{km}^{+},
\end{multline}
where we have used superscripts to distinguish the elastic tensors of the two half-spaces. These equations are to hold for all $(x_{1},x_{2})$, and the only way this can happen non-trivially is if the exponential factors are equal, and hence can be cancelled. In this manner we obtain the appropriate generalisation of Snell's law for elastic waves:
\begin{equation}
p_{i}=p_{ik}^{-}=p_{ik}^{+}, \quad i=1,2.
\end{equation}
In words, the horizontal slowness of all plane waves must be equal.

For our problem the values of $p_{i}$ and $a_{i}$ are assumed to be given, subject to the requirement that they satisfy the appropriate Christoffel equation. From the incident wave, the horizontal component of the slowness for all reflected and transmitted waves are fixed, and the problem is reduced to finding their vertical slownesses. Mathematically this requires, in general, the solution of a sixth-order polynomial equation, and hence must be done numerically. Geometrically, however, the problem is one of finding the intersections of the vertical line in slowness space fixed by $(p_{1},p_{2})$ with the slowness surfaces for the two half spaces. For the reflected waves we seek solutions with negative vertical slowness, while for the transmitted ones we need a positive vertical slowness. Two examples are shown in Fig.2. In the first case there are six real roots, with three corresponding to reflected waves, and three to transmitted waves. To each root there is a corresponding polarisation vector that can be calculated up to their amplitude. Writing, for example, $a_{im}^{-}=c_{m}^{-}\hat{a}_{im}^{-}$ with $c_{m}^{-}$ the unknown amplitude, the continuity conditions reduce to
\begin{align}
a_{i}+\sum_{m}a_{im}^{-}c_{m}^{-}&=\sum_{m}\hat{a}_{im}^{+}c_{m}^{+} \\
p_{l}A_{i3kl}^{-}a_{k}+\sum_{m}p_{lm}^{-}A_{i3kl}^{-}a_{km}^{-}c_{m}^{-}&=\sum_{m}p_{lm}^{+}A_{i3kl}^{+}\hat{a}_{km}^{+}c_{m}^{+}.
\end{align}
Here we have six simultaneous linear equations for the six unknown amplitudes of the reflected and transmitted waves. Assuming that the equations are linearly independent, they can then be solved, and in this manner we obtain the appropriate reflection and transmission coefficients for the problem. Note, in particular, that a single incident plane wave will generally excite three reflected and three transmitted waves. The excitation of planes waves of a different types by an incident wave is known as a wave conversion, and this process accounts for much of the complexity in observed seismograms.

In the second example shown in Fig.2 we can see only two roots corresponding to reflected waves, and one for transmitted waves. This gives us only three linearly independent terms to try to match the six continuity conditions, and in general this will not be possible. Here, therefore, we are required to look for complex values of the vertical slownesses. Because we are solving a sixth-order polynomial there will always be six intersections with each slowness surface, and we assume again that there are three corresponding to down-going reflected waves, and three for up-going transmitted waves. Note that for down-going waves there is a requirement on the sign of the imaginary part of the vertical slowness, and similarly for up-going waves. By including such evanescent plane waves, with amplitudes that decay exponentially away from the boundary, we can obtain sufficiently many terms to satisfy the continuity conditions, and hence again arrive at suitable reflection and transmission coefficients.

As a final comment, the above calculations depended on the linear dependence of the polarisation vectors and associated tractions for the reflected and transmitted waves. You might ask if this will always hold, and the answer is no. In such cases there exist non-trivial solutions to the equations of motion and all boundary conditions that do not need to be excited. This is a kind of eigenvalue problem that has non-trivial solutions only for certain values of frequency and slowness. When these ideas are applied to the free surface, it leads to what are known as surface waves, while at an interface they are known as Stoneley waves. There is not time to enter into this topic further\footnote{In past years surface waves were considered in further detail within the course, and so you might see reference to them in older past paper questions. These can be safely ignored.}.

\section*{What you need to know and be able to do}
\begin{itemize}
    \item[(i)] How to derive the Christoffel equation from the linearised equations of motion.
    \item[(ii)] Reduction of the problem to a symmetric eigenvalue problem, and a qualitative description of its solution.
    \item[(iii)] How to solve the Christoffel equation for an isotropic solid, and to deduce the properties of plane P- and S-waves.
    \item[(iv)] How to apply first-order perturbation theory to study the effect of slight anisotropy on plane waves.
    \item[(v)] Describe qualitatively the reflection, transmission, and conversion of elastic waves at a plane boundary. You will not be asked to perform any detailed calculations of reflection and transmission coefficients.
\end{itemize}

\end{document}
